% ID: SEL-D0044-OA-E03
% ESTADO: validada. S045, cotejada con examen y criterios.
Tomamos las dos elecciones como independientes, tal como
se modela al seleccionar dos adultos de una población muy grande.
Sean $N$ no fumar y $O$ fumar ocasionalmente.

\textbf{a)} La probabilidad de dos no fumadores es
$P(N)^2=(0.65)^2=\boxed{0.4225}$.

\textbf{b)} Hay dos órdenes disjuntos, $NO$ y $ON$.
Cada uno tiene probabilidad $0.65(0.15)$, de modo que el
resultado es $2(0.65)(0.15)=\boxed{0.195}$.
